\chapter{Reconstruction and Classification}
\label{ch:classification}

\section{Generators, relations, and the nonzero construction}

Throughout this chapter \(A=(a_{ij})\) is the Cartan matrix of
a finite reduced root system, with
\(a_{ij}=\langle\alpha_j,\alpha_i^\vee\rangle\).
We establish the existence theorem before using its presentation
to prove uniqueness. Thus the reconstruction result labelled F01
below is a consequence of the full construction, rather than an
assumption used to prove it.

\begin{proposition}[The Serre relations and the triangular construction]
\label{res:F02}
In a semisimple algebra with simple-root triples the elements
\(e_i,f_i,h_i\) generate the algebra and satisfy
\[
\begin{gathered}
[h_i,h_j]=0,\quad [h_i,e_j]=a_{ij}e_j,\quad
[h_i,f_j]=-a_{ij}f_j,\quad [e_i,f_j]=\delta_{ij}h_i,\\
(\operatorname{ad}e_i)^{1-a_{ij}}e_j=0,\qquad
(\operatorname{ad}f_i)^{1-a_{ij}}f_j=0\quad(i\ne j).
\end{gathered}\tag{6.1}\label{eq:serre-relations}
\]
Conversely form the free Lie algebra on these symbols and impose
these relations. The resulting algebra \(\mathfrak g(A)\)
has a vector space decomposition
\[
 \mathfrak g(A)=\mathfrak n^-\oplus\mathfrak h\oplus\mathfrak n^+,
\]
where \(\mathfrak h\) has basis \(h_i\), and
\(\mathfrak n^\pm\) have precisely their respective positive
or negative Serre relations. In particular no \(e_i,f_i,h_i\)
has been killed by the presentation.
\end{proposition}
\begin{proof}
In an existing algebra, a positive root \(\beta\) of height
greater than one has \(\beta-\alpha_i>0\) for some \(i\).
The nonzero root-string bracket
\([e_i,\mathfrak g_{\beta-\alpha_i}]=\mathfrak g_\beta\)
gives generation by induction. Negative roots work the same way,
and the \(h_i\) form a basis of \(\mathfrak h\).
For \(i\ne j\), \(f_j\) is a highest vector of weight
\(-a_{ij}\) for the \(i\)-th triple. The
\(\mathfrak{sl}_2\) calculation gives the second Serre relation;
the first follows by reversing signs. The other relations
are the root-space bracket calculations.

We now justify carefully that the abstract presentation is
nonzero. Initially omit the Serre relations and call the algebra
\(\widetilde{\mathfrak g}\). Repeated use of
\([e_i,f_j]=\delta_{ij}h_i\), followed by Jacobi, moves any
mixed bracket into a sum of brackets using only \(e\)'s,
only \(f\)'s, or one \(h\). For example apply
\([[u,v],w]=[u,[v,w]]-[v,[u,w]]\) to reduce the number
of mixed letters adjacent to an innermost bracket.
The root-lattice grading separates positive degree, zero,
and negative degree. Hence
\(\widetilde{\mathfrak g}=
\widetilde{\mathfrak n}^-\oplus\widetilde{\mathfrak h}
\oplus\widetilde{\mathfrak n}^+\) once freeness and
independence in these pieces are verified.

For that verification let \(F\) be the free Lie algebra on
symbols \(f'_i\), and act on
\[
 V=T(f'_1,\ldots,f'_r)\otimes\mathbb C[H_1,\ldots,H_r].
\]
If a word \(w\) has weight \(-\gamma\), put
\[
 F_i(wP)=f'_iwP,\qquad
 H_i^{\rm op}(wP)=w(H_i-\gamma(h_i))P.
\]
Define \(E_i(1P)=0\) and recursively
\[
 E_i(f'_j wP)=f'_jE_i(wP)+\delta_{ij}H_i^{\rm op}(wP).
 \tag{6.2}
\]
The recursion gives \([E_i,F_j]=\delta_{ij}H_i^{\rm op}\).
The degree of \(F_j\) changes the weight by \(-\alpha_j\),
so \([H_i^{\rm op},F_j]=-a_{ij}F_j\).
Every summand of \(E_jwP\) has weight
\(-\gamma+\alpha_j\); comparing the polynomial factors
therefore gives \([H_i^{\rm op},E_j]=a_{ij}E_j\).
The \(H_i^{\rm op}\) commute. These are all the imposed
relations of \(\widetilde{\mathfrak g}\), so this is an action.
Applying a bracket of the \(F_i\)'s to \(1\) gives the
identical free-Lie bracket inside \(T(f')=U(F)\);
PBW embeds \(F\) there. Thus the negative piece is free.
Also \(H_i^{\rm op}(1)=H_i\) proves independence of the \(h_i\).
Interchanging \(e,f\) and replacing \(h\) by \(-h\)
proves the same freeness for the positive piece.

Put \(S^+_{ij}=(\operatorname{ad}e_i)^{1-a_{ij}}e_j\).
For \(k\ne i,j\), \([f_k,S^+_{ij}]=0\) follows at once.
For \(k=i\), the usual iterated \(\mathfrak{sl}_2\)
commutator has scalar
\((1-a_{ij})(a_{ij}+1-a_{ij}-1)=0\).
For \(k=j\), it is a multiple of
\((\operatorname{ad}e_i)^{1-a_{ij}}h_j\).
If \(a_{ij}=0\), this vanishes since \(a_{ji}=0\);
otherwise the exponent is at least two and the second
bracket of \(e_i\) with \(h_j\) is already zero.
Thus all \(f_k\) annihilate \(S^+_{ij}\).
The ideal \(I^+\) generated by the \(S^+_{ij}\) within the
free positive piece is stable under all \(f_k\): apply Jacobi
inductively to a bracket of positive generators with an
\(S^+_{ij}\). The terms involving a bracket with \(f_k\)
have fewer positive generators; Cartan terms only multiply
homogeneous vectors by scalars. It is also stable under
\(\mathfrak h\). The corresponding statement holds for
the negative ideal \(I^-\).
Therefore the full ideal of Serre relations is exactly
\(I^-\oplus I^+\). It has no Cartan component and no
component of degree \(\pm\alpha_i\).
Quotienting gives precisely the decomposition asserted.
\end{proof}

\begin{theorem}[Finite type and simplicity of the Serre algebra]
\label{res:F03}
The algebra \(\mathfrak g(A)\) is finite dimensional. Its nonzero
weights are exactly the prescribed roots, each with multiplicity
one, and its zero space is the \(r\)-dimensional \(\mathfrak h\).
For connected \(A\) it is simple. For arbitrary \(A\) it is
the direct sum of the simple algebras of the connected components.
\end{theorem}
\begin{proof}
Fix \(i\). The \(i\)-th triple acts locally finitely on every
generator: on itself it is the adjoint three-dimensional module;
on \(e_j\) or \(f_j\), \(j\ne i\), the Serre relation bounds
the raising or lowering string; a Cartan generator belongs
to the sum of a trivial module and that adjoint module.
If \(x,y\) lie in finite-dimensional invariant spaces, their
brackets lie in the image of the tensor product of those
spaces. Induction on bracket length shows that every element
is in a finite-dimensional module for the \(i\)-th triple.
Consequently the finite exponentials
\[
 n_i=\exp(\operatorname{ad}e_i)
       \exp(-\operatorname{ad}f_i)
       \exp(\operatorname{ad}e_i)
\]
are well-defined automorphisms. Their effect on \(\mathfrak h\)
is \(h\mapsto h-\alpha_i(h)h_i\), so they take the weight
space of \(\gamma\) to that of \(s_i\gamma\).

The triangular decomposition says that a nonzero weight
belongs to \(Q^+\) or \(-Q^+\), and that the zero space is
exactly \(\mathfrak h\). Each individual weight space is
finite dimensional, since only finitely many words have
a specified nonnegative multidegree. Suppose a positive
weight \(\gamma=\sum k_j\alpha_j\) occurs.
If its support has one element, it is a simple root:
a free Lie algebra on one generator has no degree greater
than one. Otherwise positivity of the Euclidean form gives
an \(i\) with \(\langle\gamma,\alpha_i^\vee\rangle>0\).
The weight \(s_i\gamma\) occurs and has lower height.
At least one coefficient other than the \(i\)-th remains
positive; hence it cannot be in \(-Q^+\).
It must be in \(Q^+\), and induction on height takes it
to a simple root. Thus \(\gamma\) is a root of the
given root system. The negative case is identical.
Conversely the \(n_i\)'s and the nonzero simple generators
produce all \(W\Pi=R\).
The weight spaces of simple roots are one dimensional;
the reflection isomorphisms make all root spaces
one dimensional. It follows that
\(\dim\mathfrak g(A)=r+|R|\).

Suppose first that the diagram is connected and
\(0\ne I\triangleleft\mathfrak g(A)\).
The commuting diagonal actions of the \(h_i\) decompose \(I\)
into weight spaces. If \(I\) has only a nonzero Cartan
element \(h\), some \(\alpha_i(h)\ne0\), since \(A\)
is invertible; hence \([h,e_i]\in I\) supplies a nonzero
root vector. A root vector in \(I\) can be reflected by
the \(n_i\)'s, which preserve ideals, to some \(e_j\).
Then \(h_j=[e_j,f_j]\in I\), and
\([h_j,f_j]=-2f_j\in I\). Whenever \(a_{jk}\ne0\),
\([h_j,e_k]=a_{jk}e_k\) puts \(e_k\) in \(I\).
Connectedness propagates this to every generator.
Thus \(I=\mathfrak g(A)\), proving simplicity.
For disconnected diagrams, generators in different components
commute by the relations, and the independent Cartan and
root decomposition makes their sum direct.
\end{proof}

\begin{theorem}[Reconstruction and the isomorphism theorem]\label{res:F01}
An isomorphism of the root systems of two complex semisimple
Lie algebras extends, after choosing compatible simple-root
triples, to a Lie algebra isomorphism. In particular the root
system determines the algebra.
\end{theorem}
\begin{proof}
For an algebra with root system \(R\), the generation and
relations in Proposition~\ref{res:F02} give a surjective map
\(\mathfrak g(A)\to\mathfrak g\) sending the indicated
generators to the selected triples. Both sides have dimension
\(r+|R|\), by Theorem~\ref{res:F03} and the root decomposition.
Thus this map is an isomorphism.
Apply the same construction to the other algebra, using
the root-system isomorphism to identify the ordered simple
roots and hence their Cartan matrix. Composing the two
isomorphisms gives the required extension.
The choices of simple root vectors are compatible because
each negative vector can be normalized to have bracket
equal to its coroot.
\end{proof}

\begin{corollary}[Complex semisimple classification]\label{res:F04}
The complex simple Lie algebras are the algebras of types
\[
 A_n,\ B_n,\ C_n,\ D_n,\ E_6,E_7,E_8,F_4,G_2
\]
with the rank restrictions and repetitions specified in
Theorem~\ref{res:E05}. Every complex semisimple Lie algebra
is a direct sum of these, uniquely up to permutation.
\end{corollary}
\begin{proof}
Every semisimple algebra gives a root system by
Proposition~\ref{res:D05}; its irreducible components
are the finite list in Theorem~\ref{res:E05}.
The coordinate models and Serre construction prove that
each component actually occurs, and reconstruction proves
uniqueness. The simple-ideal decomposition
of Theorem~\ref{res:B06} agrees with the root-system
decomposition: different ideals have orthogonal Killing
forms and disjoint root systems, and an irreducible
component yields a simple Serre algebra.
For \(A,B,C,D\), Example~\ref{res:D08} identifies these
algebras with the stated classical matrix algebras.
\end{proof}

\section{Integral structure constants}

\begin{theorem}[Chevalley bases]\label{res:F05}
There are nonzero root vectors \(x_\alpha\) such that
\[
 [x_\alpha,x_{-\alpha}]=h_\alpha,\qquad
 [x_\alpha,x_\beta]=N_{\alpha\beta}x_{\alpha+\beta},
\]
and
\[
 N_{-\alpha,-\beta}=-N_{\alpha\beta},\qquad
 N_{\alpha\beta}=\pm(p+1)
\]
when \(\alpha+\beta\) is a root and the \(\alpha\)-string
through \(\beta\) begins at \(\beta-p\alpha\).
Together with the simple coroots these vectors are a basis
whose structure constants are integers.
\end{theorem}
\begin{proof}
The Serre relations admit the involution
\[
 \sigma(e_i)=-f_i,\quad \sigma(f_i)=-e_i,\quad
 \sigma(h_i)=-h_i.
\]
Substitution in (6.1) verifies every relation, and its square
is the identity. Hence it carries
\(\mathfrak g_\alpha\) to \(\mathfrak g_{-\alpha}\).
For each pair \(\{\alpha,-\alpha\}\), choose \(x_\alpha\)
and set \(x_{-\alpha}=-\sigma x_\alpha\).
Opposite root spaces pair perfectly, so their bracket is
a nonzero multiple of \(h_\alpha\). Multiplying both vectors
by a suitable complex scalar changes that multiple by its
square; choose it to obtain bracket \(h_\alpha\).
Applying \(\sigma\) to a bracket immediately gives
\(N_{-\alpha,-\beta}=-N_{\alpha\beta}\).

We verify the magnitude, since a sign choice alone would
not prove integrality. In a root string
\(\beta-p\alpha,\ldots,\beta+q\alpha\),
the \(\mathfrak{sl}_2\) calculation gives
\[
 [x_{-\alpha},[x_\alpha,x_\beta]]
       =q(p+1)x_\beta .
 \tag{6.3}
\]
Apply Jacobi twice to
\([[x_\alpha,x_\beta],[x_{-\alpha},x_{-\beta}]]\).
Using (6.3) for both strings, and the negative-constant
identity, gives
\[
 N_{\alpha\beta}^2h_{\alpha+\beta}
       =q'(p'+1)h_\alpha+q(p+1)h_\beta,
 \tag{6.4}
\]
where \(p',q'\) describe the \(\beta\)-string through
\(\alpha\). Compare the coefficient of \(t_\beta\),
using \(h_\gamma=2t_\gamma/(\gamma,\gamma)\):
\[
 N_{\alpha\beta}^2
   =q(p+1)\frac{(\alpha+\beta,\alpha+\beta)}{(\beta,\beta)}.
 \tag{6.5}
\]
The rank two lists give the following complete possibilities
when \(q>0\):
\[
\begin{array}{c|rrrrrr}
(p,q)&(0,1)&(0,2)&(0,3)&(1,1)&(1,2)&(2,1)\\
(\alpha+\beta)^2/\beta^2&1&1/2&1/3&2&1&3 .
\end{array}
\]
In each case the right side of (6.5) is \((p+1)^2\).
Thus \(N_{\alpha\beta}=\pm(p+1)\).
If \(\alpha+\beta\) is neither zero nor a root the bracket
is zero; opposite roots have the normalized coroot bracket.
Finally every coroot is an integral combination of the simple
coroots, by applying the simple-root theorem to the dual root
system. The remaining Cartan brackets have their integral
Cartan entries. All structure constants are therefore integers.

The construction chooses all signs simultaneously through
one involution; it does not attempt to assign pairwise signs
independently. Under \(x_\alpha\mapsto c_\alpha x_\alpha\)
the opposite normalization requires \(c_\alpha c_{-\alpha}=1\),
and a bracket constant changes by
\(c_\alpha c_\beta/c_{\alpha+\beta}\).
These equations specify precisely the comparison of normalizations.
\end{proof}

\section{Borel subalgebras without an omitted fixed-point argument}

We include the two analytic ingredients used below, so that
the conjugacy theorem does not depend on a later classification
of compact groups.

\begin{lemma}[A matrix-group fact]\label{supp:closed-matrix-groups}
A closed subgroup \(H\subset\operatorname{GL}(V_{\mathbb R})\)
is a Lie group with Lie algebra
\(\mathfrak h=\{X:\exp(tX)\in H\text{ for all real }t\}\).
Its identity component is generated by exponentials from
\(\mathfrak h\).
\end{lemma}
\begin{proof}
The product limits
\[
 (e^{tX/n}e^{tY/n})^n\longrightarrow e^{t(X+Y)},\qquad
 (e^{sX}e^{sY}e^{-sX}e^{-sY})^{n}
       \longrightarrow e^{t[X,Y]},\quad s=\sqrt{t/n},
\]
first for \(t\ge0\) and then using inverses, follow by
expanding the matrix products through orders \(1/n\).
They prove that \(\mathfrak h\) is a vector space and
closed under brackets; closure of \(H\) is used in passing
to the limits. It is also closed as a set.
Choose a vector-space complement \(M\) to \(\mathfrak h\).
The map \((X,Y)\mapsto e^Xe^Y\) from
\(\mathfrak h\oplus M\) has invertible differential at zero,
and hence is a local coordinate system near the identity.
If there were arbitrarily small nonzero \(Y_j\in M\)
with \(e^{Y_j}\in H\), a subsequence of
\(Y_j/\|Y_j\|\) would approach a unit \(Y\in M\).
Powers with exponent \(\lfloor t/\|Y_j\|\rfloor\)
would give \(e^{tY}\in H\) for every \(t\), a contradiction.
Thus in a sufficiently small coordinate neighborhood an
element \(e^Xe^Y\in H\) has \(Y=0\).
This gives the submanifold chart \(\exp(\mathfrak h)\)
and the asserted tangent space. Translates give the global
Lie structure. The subgroup generated by these local
exponentials is open in \(H\); its cosets are also open,
so it is exactly the identity component.
\end{proof}

\begin{lemma}[The compact form used for the flag orbit]
\label{supp:compact-form-flag}
There is a conjugate-linear involution \(\tau\) of
\(\mathfrak g\) whose fixed algebra \(\mathfrak u\)
has negative definite Killing form. For each positive system
and \(\mathfrak b=\mathfrak h\oplus\mathfrak n^+\),
\[
 \mathfrak g_{\mathbb R}=\mathfrak u+\mathfrak b_{\mathbb R}.
\]
The inner orbit of \(\mathfrak b\) in the Grassmannian of
\(\dim\mathfrak b\) dimensional subspaces is compact.
\end{lemma}
\begin{proof}
Use the real span of the Chevalley basis and define
\(\tau(h_i)=-h_i\), \(\tau(x_\alpha)=-x_{-\alpha}\),
conjugating scalars. The real structure constants and their
negative-constant identity prove that \(\tau\) preserves
brackets and has square one.
Its fixed real basis consists of
\[
 ih_i,\qquad x_\alpha-x_{-\alpha},\qquad
 i(x_\alpha+x_{-\alpha})\quad(\alpha>0).
 \tag{6.6}
\]
On the real coroot space \(B\) is positive definite,
and \(B(x_\alpha,x_{-\alpha})=2/(\alpha,\alpha)>0\).
Orthogonality of root spaces makes (6.6) a negative
definite decomposition. In particular
\((v,w)_{\rm H}=-B(v,\tau w)\) is a positive definite
Hermitian form on \(\mathfrak g\).
The displayed real sum follows by writing a negative root
vector as the sum of an element from (6.6) and a positive
root vector, and by splitting Cartan coefficients into
real and imaginary parts.

Let \(G=\operatorname{Aut}(\mathfrak g)^0\).
It is a closed matrix group: preserving each bracket is
a finite set of polynomial equations in the matrix entries.
The matrix-group lemma applies, and its Lie algebra is
\(\operatorname{Der}(\mathfrak g)=\operatorname{ad}\mathfrak g\),
by Theorem~\ref{res:B07}. Thus it is the inner group used
in Chapter~4.
Let \(K\) be the identity component of the automorphisms
preserving the positive Hermitian form. It is compact.
Its Lie algebra is \(\operatorname{ad}\mathfrak u\):
the adjoint of \(\operatorname{ad}z\) is
\(-\operatorname{ad}(\tau z)\), so skew adjointness is
equivalent, by the zero center, to \(\tau z=z\).
The tangent to the \(G\)-orbit at \(\mathfrak b\) is the
image of \(\mathfrak g_{\mathbb R}\) in the tangent
space of the Grassmannian, with kernel its real normalizer.
The sum \(\mathfrak g_{\mathbb R}=\mathfrak u+
\mathfrak b_{\mathbb R}\) shows that \(K\)'s orbit has
the same tangent space. The inverse function theorem
therefore makes this compact orbit open in the \(G\)-orbit.
It is closed there by compactness. The \(G\)-orbit is
connected, so both orbits are equal. This proves compactness.
\end{proof}

\begin{lemma}[Solvable fixed points on a compact projective set]
\label{supp:solvable-projective-fixed}
Let a complex solvable Lie algebra \(\mathfrak a\) act on a
finite-dimensional space \(V\). Every nonempty compact subset
\(X\subset\mathbb P(V)\), invariant under
\(\exp(zx)\) for every \(x\in\mathfrak a,z\in\mathbb C\),
contains a point fixed by \(\mathfrak a\).
\end{lemma}
\begin{proof}
Induct on \(\dim\mathfrak a\). Choose a codimension one ideal
\(\mathfrak a'\), which exists because a nonzero solvable
algebra has a proper derived algebra.
By induction its fixed-point subset \(Y\subset X\)
is nonempty. It is closed, hence compact, and is invariant
under \(\mathfrak a\): conjugation by \(e^{zx}\) preserves
\(\mathfrak a'\). Choose \(a\) projecting to a generator
of \(\mathfrak a/\mathfrak a'\) and a point \([v]\in Y\).
Multiply \(a\) by a nonzero complex scalar so that its
distinct eigenvalues have distinct real parts.
In the Jordan decomposition of \(v\), take the eigenvalue
of largest real part present, and, within it, the highest
nonzero nilpotent power. As \(t\to+\infty\), the line
\([e^{ta}v]\) tends to the line of that highest-power
vector. This is an eigenline: lower real parts disappear
exponentially and lower powers disappear after division
by the leading power of \(t\). The limit belongs to \(Y\)
and is fixed by \(a\). It was already fixed by
\(\mathfrak a'\), so is fixed by \(\mathfrak a\).
\end{proof}

\begin{theorem}[Borel subalgebras]\label{res:F06}
The algebra \(\mathfrak b=\mathfrak h\oplus
\bigoplus_{\alpha>0}\mathfrak g_\alpha\) is maximal solvable.
Every solvable subalgebra lies in an inner conjugate of
\(\mathfrak b\); consequently all maximal solvable
subalgebras, called Borel subalgebras, are conjugate.
Moreover \(N_{\mathfrak g}(\mathfrak b)=\mathfrak b\).
\end{theorem}
\begin{proof}
Brackets increase positive height, so \(\mathfrak n^+\)
is nilpotent. The quotient \(\mathfrak b/\mathfrak n^+\)
is abelian; thus \(\mathfrak b\) is solvable.
A subalgebra strictly containing it is \(\mathfrak h\)-stable,
and hence contains some negative root space.
It then contains that root's full \(\mathfrak{sl}_2\),
so cannot be solvable. The same weight argument shows that
an element normalizing \(\mathfrak b\) has no negative
root component, proving its stated normalizer.

Embed the Grassmannian into
\(\mathbb P(\bigwedge^d\mathfrak g)\), \(d=\dim\mathfrak b\),
by \(L\mapsto[\bigwedge^d L]\).
The inner orbit \(X\) of \(\mathfrak b\) is nonempty and
compact by Lemma~\ref{supp:compact-form-flag}.
A solvable subalgebra \(\mathfrak a\) acts on this
exterior power, preserves \(X\), and the preceding fixed-point
lemma supplies an \(\mathfrak a\)-fixed line belonging to
some conjugate \(\mathfrak b'\).
A line \(\bigwedge^d\mathfrak b'\) is fixed precisely when
\([\mathfrak a,\mathfrak b']\subset\mathfrak b'\):
the infinitesimal action modulo \(\mathfrak b'\) is
the corresponding tangent map of the Grassmannian.
Thus \(\mathfrak a\subset N_{\mathfrak g}(\mathfrak b')
=\mathfrak b'\). If \(\mathfrak a\) was maximal solvable,
this inclusion is equality.
\end{proof}

\begin{theorem}[Parabolic subalgebras]\label{res:F07}
Subalgebras containing \(\mathfrak b\) are in bijection with
subsets \(S\subset\Pi\):
\[
 \mathfrak p_S=
 \mathfrak b\oplus\bigoplus_{\alpha\in R_S^+}
                       \mathfrak g_{-\alpha},
 \qquad R_S=R\cap\operatorname{span}_{\mathbb R}S .
\]
They are the standard parabolic subalgebras.
Every subalgebra containing a Borel is inner conjugate
to one of these. There is a decomposition
\[
 \mathfrak p_S=\mathfrak l_S\ltimes\mathfrak u_S,\quad
 \mathfrak l_S=\mathfrak h\oplus
                 \bigoplus_{\alpha\in R_S}\mathfrak g_\alpha,
 \quad
 \mathfrak u_S=\bigoplus_{\alpha\in R^+\setminus R_S}
                                       \mathfrak g_\alpha.
\]
Here \(\mathfrak l_S\) is reductive,
\(\mathfrak u_S\) is the nilpotent radical,
and replacing positive by negative gives the opposite
parabolic with the same Levi subalgebra.
\end{theorem}
\begin{proof}
Root addition verifies that \(\mathfrak p_S\) is a subalgebra
and \(\mathfrak u_S\) an ideal; height makes the latter
nilpotent. The generators from \(S\), their negatives,
and their coroots give the semisimple Serre algebra
for that root subsystem. Its complementary Cartan space
\(\{h:\alpha(h)=0\ \forall\alpha\in S\}\)
is central in \(\mathfrak l_S\). Thus \(\mathfrak l_S\)
is their direct sum and is reductive.

Conversely let \(\mathfrak p\supset\mathfrak b\).
It is a sum of \(\mathfrak h\) and root spaces, and define
\(S=\{\alpha_i:\mathfrak g_{-\alpha_i}\subset\mathfrak p\}\).
If \(\mathfrak g_{-\beta}\subset\mathfrak p\) for a nonsimple
positive root \(\beta\), choose \(\alpha_i\) with
\(\gamma=\beta-\alpha_i>0\) a root.
The two nonzero brackets
\[
 [\mathfrak g_{\alpha_i},\mathfrak g_{-\beta}]
       =\mathfrak g_{-\gamma},\qquad
 [\mathfrak g_\gamma,\mathfrak g_{-\beta}]
       =\mathfrak g_{-\alpha_i}
\]
belong to \(\mathfrak p\).
Induction on height puts every simple root in the support
of \(\gamma\), and also \(\alpha_i\), in \(S\).
Thus \(\beta\in R_S\).
The converse inclusion follows by generation of the
negative subsystem from its simple roots.
This proves \(\mathfrak p=\mathfrak p_S\) and uniqueness of \(S\).
Borel conjugacy proves the assertion about arbitrary parabolics.

To verify that the displayed ideal is the entire nilpotent
radical, let \(I\) be a nilpotent ideal of \(\mathfrak p_S\).
Its image in the reductive quotient lies in
\(\mathfrak z(\mathfrak l_S)\); any ideal with a nonzero
simple-factor image is not nilpotent.
As \(I\) is \(\mathfrak h\)-stable, a nonzero zero-weight
part of that image is represented by
\(0\ne z\in I\cap\mathfrak z(\mathfrak l_S)\).
Some root outside \(R_S\) has \(\alpha(z)\ne0\),
because all roots span and the roots in \(S\) kill \(z\).
Replacing it by its positive sign, its root vector
\([z,x_\alpha]/\alpha(z)\) belongs to \(I\).
But then \(\operatorname{ad}_I z\) has a nonzero eigenvalue,
contradicting nilpotence of \(I\).
Thus its quotient image is zero and \(I\subset\mathfrak u_S\).
The root formulas with negative roots give the opposite
parabolic and the common Levi factor.
\end{proof}

The Serre proof is based on
\cite[Sections 18--19]{humphreys} and
\cite[Section 24]{etingof}, including the free negative-word
module that prevents a degenerate presentation.
The integral normalization is
\cite[Section 25]{humphreys}. The last section supplies the
fixed-point and compact-orbit arguments behind the Borel
and parabolic assertions in \cite[Sections 51.3--51.5]{etingof}.
